% Appendix A (two-cause rewrite, 2026-11): settings and baselines. Settings are not results;
% the baseline tables are generated.
\chapter{Settings and Baselines}
\label{app:supplementary}

\section{Training Settings}
\label{app:hyperparameters}

Table~\ref{tab:hyperparameters} lists the settings shared by all runs and the settings of
our method. The protocol is described in Section~\ref{sec:e-protocol}.

\begin{table}[!htbp]
  \centering
  \caption[Training settings]{\textbf{Training settings.} Top: shared by all methods.
  Bottom: the proposed method.}
  \label{tab:hyperparameters}
  \small
  \begin{tabular}{@{}l >{\raggedright\arraybackslash}p{9.6cm}@{}}
    \toprule
    Setting & Value \\
    \midrule
    \multicolumn{2}{@{}l}{\emph{Shared}} \\
    Main encoder & LayoutLMv3-base ($12$ layers, width $768$) \\
    Optimiser & AdamW, learning rate $5\times10^{-5}$, weight decay $0.01$ \\
    Gradient clipping & norm $1.0$ \\
    Batch size & $2$ documents, gradient checkpointing \\
    Precision & full precision \\
    Epochs per session & at most $10$ ($100$ for joint training) \\
    Early stopping & $10\%$ of the session held out; patience $2$ epochs; minimum gain $0.1$ F1 \\
    Input & $512$ tokens, $224\times224$ image, boxes in $[0,1000]$ \\
    Runs per cell & $3$ \\
    \midrule
    \multicolumn{2}{@{}l}{\emph{CoLaR-SBR}} \\
    Injection layer $k$ & $4$ (layers $1$ to $4$ frozen after the first session) \\
    Rank $r$ of the stored factors & $128$, half precision \\
    Stored documents & $50$ per session ($10$ on CIL-FUNSD) \\
    Replayed documents per batch & $4$ \\
    Claim threshold & $0.5$ \\
    \midrule
    \multicolumn{2}{@{}l}{\emph{SBR-Recall}} \\
    Rank of an update & $64$, scale $\alpha=128$ \\
    Updated maps & query, key, value, attention output and both feed-forward maps of layers $5$ to $12$ \\
    Stored documents & none on one schema; $50$ per session across schemas \\
    \bottomrule
  \end{tabular}
\end{table}

\section{Baselines}
\label{app:baselines}

Table~\ref{tab:baselinecard-full} gives the full card of the eight baselines: what our
re-implementation keeps from the published method, what it changes, and what it keeps
across the sessions. Table~\ref{tab:baseline-hp} lists their hyper-parameters as recorded
by the runs; they are the defaults of each method. LCA keeps its published optimiser
instead of AdamW, i.e., SGD with momentum.

\begin{landscape}
\begin{table}[p]
  \centering
  \caption[Baseline card]{\textbf{Baseline card}: the published method, what our
  re-implementation keeps and changes, and its memory on CIL-CORD.}
  \label{tab:baselinecard-full}
  \resizebox{\linewidth}{!}{% AUTO-GENERATED by scripts/rca/make_clean_tables.py -- do not edit.
\begin{tabular}{@{}lll>{\raggedright\arraybackslash}p{4.0cm}>{\raggedright\arraybackslash}p{4.6cm}>{\raggedright\arraybackslash}p{2.6cm}@{}}
\toprule
Method & Venue & Family & Kept from the method & Changed in our version & Memory \\
\midrule
ER & NeurIPS 2019 & Replay & reservoir buffer replayed in every step & -- & 250 raw documents \\
DER++ & NeurIPS 2020 & Replay & replay of documents with labels and stored logits & -- & 250 raw documents + logits \\
LwF & ECCV 2016 & Regularisation & distillation of the old logits towards the previous model & -- & previous model \\
CODA-Prompt & CVPR 2023 & Prompt, LoRA & attention-weighted prompt components with an orthogonality penalty, frozen encoder & prompts prepended to the input sequence (an input-prepend variant) & none (prompt pool) \\
CL-LoRA & CVPR 2025 & Prompt, LoRA & orthogonality penalty, distillation towards the previous model & one shared rank-8 adapter retrained each session (published: random-orthogonal shared and block-wise task-specific adapters, early-exit distillation); gradient reassignment at scale 1, i.e.\ off & previous model, past adapter matrices \\
CPFD & EMNLP 2023 & Continual NER & confidence pseudo-labels, logit and attention-map distillation, adaptive CE factor & one growing token head (published: per-task classifiers); sample weights set to 1 & previous model \\
IS3 & Findings ACL 2024 & Continual NER & distillation on ``O'' tokens, head-row gradient surgery, prototype replay & one growing token head (published: per-task classifiers) & previous model, class prototypes \\
LCA & ICLR 2026 & Merging & per-session fine-tuning (SGD), TIES backbone merge, Gaussian classifier realignment & one growing token head (published: per-task heads); Gaussians per token class, ``O'' excluded & base and per-session backbones, class Gaussians \\
\bottomrule
\end{tabular}
}
\end{table}
\end{landscape}

\begin{table}[!htbp]
  \centering
  \caption[Hyper-parameters of the baselines]{\textbf{Hyper-parameters of the baselines}
  as recorded by the runs on CIL-CORD. The last three columns give the protocol and the
  number of runs per stream (clean: the protocol of Section~\ref{sec:e-protocol}).}
  \label{tab:baseline-hp}
  \resizebox{\textwidth}{!}{% AUTO-GENERATED by scripts/rca/make_clean_tables.py -- do not edit.
\begin{tabular}{@{}l>{\raggedright\arraybackslash}p{9.6cm}ccc@{}}
\toprule
 & & \multicolumn{3}{c}{Protocol, seeds} \\
\cmidrule(l){3-5}
Method & Hyper-parameters (recorded, CIL-CORD) & CIL-CORD & Mixed & DIL \\
\midrule
ER & buffer\_size 250, replay\_batch\_size 8, epochs 10 & clean, 3 & clean, 3 & clean, 3 \\
DER++ & buffer\_size 250, alpha 0.5, beta 0.5, epochs 10 & clean, 3 & clean, 3 & clean, 3 \\
LwF & alpha 1.0, temperature 2.0, epochs 10 & clean, 3 & clean, 3 & clean, 3 \\
CODA-Prompt & n\_components 10, prompt\_length 5, lambda\_ortho 0.1, epochs 10 & clean, 3 & clean, 3 & clean, 3 \\
CL-LoRA & lora\_rank 8, lora\_alpha 16, lambda\_ortho 0.5, kd\_alpha 1.0, temperature 2.0, shared\_grad\_scale 1.0, epochs 10 & clean, 3 & clean, 3 & clean, 3 \\
CPFD & distill\_weight 2.0, threshold 0.001, epochs 10 & clean, 3 & clean, 3 & clean, 3 \\
IS3 & kd\_weight 2.0, proto\_weight 0.1, grad\_weight 0.6, epochs 10 & clean, 3 & clean, 3 & clean, 3 \\
LCA & lr 0.01, merge\_method ties, merge\_topk 100, ca\_epochs 10, ca\_samples\_per\_cls 256, epochs 10 & clean, 3 & clean, 3 & clean, 3 \\
\bottomrule
\end{tabular}
}
\end{table}

\paragraph{Choice of the baselines.} We fixed the set of baselines after a first comparison
on CIL-CORD, Mixed and DIL. EWC~\citep{kirkpatrick2017ewc}, L2P~\citep{wang2022l2p} and
DualPrompt~\citep{wang2022dualprompt} did not learn sessions~4 and~5 of CIL-CORD in any
run (Table~\ref{tab:dropped}); hence, we dropped them. CPFD and IS3 also fail on this
stream; however, we kept them because they are designed for the shift of the background
label. It is noted that the baselines differ in what they learn as well as in what they
forget: CPFD reaches $\nCpfdLateMax$ F1 on its best later session right after it is
learned, whereas IS3, CODA-Prompt and LCA each have a later session that is learned to at
most $\nBaseWorstLateMax$ F1 in every run. In contrast, naive fine-tuning learns every
session to at least $\nNaiveDiagMin$.

\begin{table}[!htbp]
  \centering
  \caption[Baselines that were dropped]{\textbf{Baselines that were dropped} (CIL-CORD,
  first comparison, earlier protocol; AA per run). Unlearned sessions: the sessions whose
  F1 right after training stays near zero.}
  \label{tab:dropped}
  \small% AUTO-GENERATED by scripts/rca/make_clean_tables.py -- do not edit.
\begin{tabular}{lcccc}
\toprule
 & \multicolumn{3}{c}{AA, seed} & Unlearned sessions \\
\cmidrule(lr){2-4}
Method & 42 & 7 & 123 & (seeds 42 / 7 / 123) \\
\midrule
Naive (reference) & 18.8 & 19.1 & 19.1 & none / none / none \\
ER (kept) & 16.9 & 14.4 & 16.6 & none / none / none \\
\midrule
EWC & 1.1 & 0.0 & 0.3 & 3--5 / 4--5 / 4--5 \\
L2P & 0.0 & 2.1 & 1.3 & 3--5 / 2--5 / 2--5 \\
DualPrompt & 0.0 & 2.0 & 3.3 & 3--5 / 2--5 / 2--5 \\
ER with SAM & 14.2 & 16.6 & 17.0 & none / none / none \\
\bottomrule
\end{tabular}

\end{table}
